\documentclass[../../main.tex]{subfiles}% 注意这里的文件路径不能用 ./main.tex ,否则用latexmk编译子文件会报错
\graphicspath{{\subfix{./image/}}{\subfix{../../image/}}} % 指定图片引用路径,后续可以直接使用图片文件名
% 如果图片存放在上级目录,则使用 ../../image/ 作为备用路径

% 例如:
% \begin{figure}[H]
% \centering
% \includegraphics[width=0.3\textwidth]{图.png}
% \caption{}
% \label{figure:图}
% \end{figure}
% 注意:上述\label{}一定要放在\caption{}之后,否则引用图片序号会只会显示??.

\begin{document}
% SOURCE p=143; printed=128; object=ch04-compact-title; region=top
\section{紧空间}\label{section:ch04-compact-spaces}
% SOURCE p=143; printed=128; object=ch04-compact-introduction; region=upper
在\cref{section:ch00-metric-spaces}中，我们给出了度量空间紧性的三种等价刻画：Heine--Borel 性质、Bolzano--Weierstrass 性质，以及完备性加全有界性。对于一般拓扑空间，只有前两种说法有意义，而第一种最为有用。因此作如下定义。
% SOURCE p=143; printed=128; object=ch04-compact-definition; region=upper
\begin{definition}[紧空间]\label{definition:ch04-compact}
拓扑空间 $X$ 称为\textbf{紧的}，如果每当 $\{U_\alpha\}_{\alpha\in A}$ 是 $X$ 的开覆盖，即每个 $U_\alpha$ 开且 $X=\bigcup_{\alpha\in A}U_\alpha$ 时，总存在有限子集 $B\subseteq A$，使 $X=\bigcup_{\alpha\in B}U_\alpha$。简言之，$X$ 紧当且仅当它的每个开覆盖都有有限子覆盖。
\end{definition}
% SOURCE p=143; printed=128; object=ch04-compact-wording; region=upper
后一种说法虽然简洁，却多少带有一词兼顾不同对象的意味：“开”修饰的是覆盖中的集合，“有限”修饰的则是集合族的成员数。
% SOURCE p=143; printed=128; object=ch04-compact-subset-definition; region=upper-middle
\begin{definition}[紧子集与预紧子集]\label{definition:ch04-compact-subset}
拓扑空间 $X$ 的子集 $Y$ 称为\textbf{紧的}，如果它在相对拓扑下是紧空间。因此，$Y\subseteq X$ 紧，当且仅当对任意满足 $Y\subseteq\bigcup_{\alpha\in A}U_\alpha$ 的 $X$ 中开集族，存在有限 $B\subseteq A$ 使 $Y\subseteq\bigcup_{\alpha\in B}U_\alpha$。若 $Y$ 的闭包紧，则称 $Y$ \textbf{预紧}。
\end{definition}
% SUPPLEMENT: id=v14-ch04-compact-definition-details; parent=ch04-compact-subset-definition; kind=proof+motivation+hypotheses; origin=independent-derivation
\begin{proof}
\textbf{相对拓扑刻画的补写（重构新增）。}若 $Y$ 相对紧，把每个环境开集 $U_\alpha$ 与 $Y$ 相交，就得到 $Y$ 的相对开覆盖，有限子覆盖对应原族的有限覆盖。反之，$Y$ 的每个相对开集 $V_\alpha$ 都可写成 $Y\cap U_\alpha$，其中 $U_\alpha$ 在 $X$ 中开；若这些 $V_\alpha$ 覆盖 $Y$，则这些 $U_\alpha$ 也覆盖 $Y$。环境开覆盖的有限化再与 $Y$ 相交，即为相对开覆盖的有限化。
\end{proof}
\begin{remark}
\textbf{动机与逐条件说明（重构新增）。}紧性让任意多项局部开条件的覆盖，能够由有限项承担。这正是后面将逐点分离合并为整体分离的机制。量词必须遍历所有开覆盖，而不只是一个选定覆盖；子覆盖必须从原覆盖中选取，不能另造一个开集 $X$ 来覆盖。若把“开覆盖”改成任意子集覆盖，所有单点组成的覆盖会迫使空间有限，便失去紧性的本意。这里不把 Hausdorff 纳入紧的定义，遇到需要分离性的结论时另列条件。

空空间紧，因为空子族已经覆盖它；任意有限空间也紧，只需为每个点选一个覆盖成员。预紧强调的是在指定环境空间中取闭包后的性质，因而依赖该环境。例如 $(0,1)$ 在 $\mathbb R$ 中预紧但不紧；开覆盖 $\{(1/n,1):n\ge2\}$ 没有有限子覆盖。该集合若以自身为环境，闭包就是自身，因而不预紧。“预紧”在本节的定义是闭包紧；不能在一般拓扑空间中直接换成需要度量的“全有界”。
\end{remark}
% END SUPPLEMENT: id=v14-ch04-compact-definition-details
% SOURCE p=143; printed=128; object=ch04-fip-introduction; region=middle
De Morgan 律给出紧性的如下闭集刻画。
% SOURCE p=143; printed=128; object=ch04-finite-intersection-property; region=middle
\begin{definition}[有限交性质]\label{definition:ch04-finite-intersection-property}
$X$ 的子集族 $\{F_\alpha\}_{\alpha\in A}$ 称为具有\textbf{有限交性质}，如果对每个有限 $B\subseteq A$ 都有 $\bigcap_{\alpha\in B}F_\alpha\ne\varnothing$。
\end{definition}
% SUPPLEMENT: id=v14-ch04-fip-remark; parent=ch04-finite-intersection-property; kind=motivation+hypotheses; origin=independent-derivation
\begin{remark}
\textbf{动机与逐条件说明（重构新增）。}有限交性质表示任意有限多项约束可以同时满足；紧性将把闭约束的这种有限相容性提升为全部相容。必须检查任意有限多个集合，两两相交不足以替代：在 $\{1,2,3\}$ 上，$\{1,2\},\{2,3\},\{1,3\}$ 两两相交，但三者交为空。采用空交等于 $X$ 的约定，所以有限交性质也检查 $B=\varnothing$；当 $X=\varnothing$ 时，没有任何族满足此性质。这一约定使后面的等价刻画也适用于空空间。
\end{remark}
% END SUPPLEMENT: id=v14-ch04-fip-remark
% SOURCE p=143; printed=128; object=ch04-compact-fip; number=4.21; region=middle
\begin{proposition}\label{proposition:ch04-compact-fip}
拓扑空间 $X$ 紧，当且仅当每个具有有限交性质的闭集族 $\{F_\alpha\}_{\alpha\in A}$ 都满足 $\bigcap_{\alpha\in A}F_\alpha\ne\varnothing$。
\end{proposition}
% SOURCE p=143; printed=128; object=ch04-proof-compact-fip; region=middle
\begin{proof}
令 $U_\alpha=F_\alpha^c$。则每个 $U_\alpha$ 开，且 $\bigcap_{\alpha\in A}F_\alpha\ne\varnothing$ 当且仅当 $\bigcup_{\alpha\in A}U_\alpha\ne X$；又，$\{F_\alpha\}$ 有有限交性质，当且仅当 $\{U_\alpha\}$ 没有任何有限子族覆盖 $X$。结论随即得到。
\end{proof}
% SUPPLEMENT: id=v14-ch04-compact-fip-details; parent=ch04-compact-fip; kind=proof-detail+motivation+hypotheses; origin=independent-derivation
\begin{proof}
\textbf{两个方向的补写（重构新增）。}若 $X$ 紧而上述闭集总交为空，其补集便是 $X$ 的开覆盖。紧性给有限补集覆盖，取补得相应有限闭集的交为空，违反有限交性质。反之，若闭集族的性质成立，而某开覆盖没有有限子覆盖，则其补集族有有限交性质，却因原族覆盖 $X$ 而总交为空，矛盾。故任意开覆盖都有有限子覆盖。
\end{proof}
\begin{remark}
\textbf{动机与条件说明（重构新增）。}闭性让补集成为开集，是把定义转换成覆盖问题的关键，不能任意去掉。即使 $X=[0,1]$ 紧，子集族 $F_n=(0,1/n)$ 也有有限交性质而总交为空，但这些集合不闭。不要求指标族可数；只检查可数闭集族得到的是后面较弱的可数紧性。
\end{remark}
% END SUPPLEMENT: id=v14-ch04-compact-fip-details
% SOURCE p=143; printed=128; object=ch04-compact-facts-intro; region=middle-lower
现在列出紧空间的几个基本事实。
% SOURCE p=143; printed=128; object=ch04-closed-subset-compact; number=4.22; region=lower
\begin{proposition}\label{proposition:ch04-closed-subset-compact}
紧空间的闭子集是紧的。
\end{proposition}
% SOURCE p=143; printed=128; object=ch04-proof-closed-subset-compact; region=lower
\begin{proof}
设 $X$ 紧、$F\subseteq X$ 闭，且 $\{U_\alpha\}_{\alpha\in A}$ 是 $X$ 中覆盖 $F$ 的开集族。则 $\{U_\alpha\}_{\alpha\in A}\cup\{F^c\}$ 是 $X$ 的开覆盖，因而有有限子覆盖。若其中含 $F^c$，将它去掉，剩下的有限个 $U_\alpha$ 仍覆盖 $F$。
\end{proof}
% SUPPLEMENT: id=v14-ch04-closed-compact-remark; parent=ch04-closed-subset-compact; kind=proof-detail+motivation+hypotheses; origin=independent-derivation
\begin{remark}
\textbf{动机与逐条件说明（重构新增）。}闭性使补集 $F^c$ 开，从而可以把只覆盖 $F$ 的局部问题补成覆盖整个 $X$ 的问题；去掉 $F^c$ 不会丢失 $F$ 中任何点，因为 $F\cap F^c=\varnothing$。若 $F=\varnothing$，剩下的空族也允许。这里不需要 Hausdorff。闭性不能普遍删除，例如紧空间 $[0,1]$ 的子集 $(0,1)$ 不紧；环境紧性也不能删除，例如 $\mathbb R$ 在自身中闭而不紧，开覆盖 $\{(-n,n):n\in\mathbb N\}$ 没有有限子覆盖。
\end{remark}
% END SUPPLEMENT: id=v14-ch04-closed-compact-remark
% SOURCE p=143; printed=128; object=ch04-point-compact-separation; number=4.23; region=lower
\begin{proposition}\label{proposition:ch04-point-compact-separation}
若 $F$ 是 Hausdorff 空间 $X$ 的紧子集，且 $x\notin F$，则存在不交开集 $U,V$，使 $x\in U$ 且 $F\subseteq V$。
\end{proposition}
% SOURCE p=143; printed=128; object=ch04-proof-point-compact-separation; region=lower
\begin{proof}
对每个 $y\in F$，选不交开集 $U_y,V_y$，使 $x\in U_y$、$y\in V_y$。族 $\{V_y\}_{y\in F}$ 是 $F$ 的开覆盖，故存在有限子覆盖 $\{V_{y_j}\}_{j=1}^n$。于是 $U=\bigcap_{j=1}^nU_{y_j}$、$V=\bigcup_{j=1}^nV_{y_j}$ 具有所需性质。
\end{proof}
% SUPPLEMENT: id=v14-ch04-point-compact-details; parent=ch04-point-compact-separation; kind=proof-detail+motivation+hypotheses; origin=independent-derivation
\begin{remark}
\textbf{动机、条件与边界核对（重构新增）。}Hausdorff 性及 $x\notin F$ 使每对不同点 $x,y$ 可分离；紧性把无限多个局部分离压缩成有限多个，因此 $U$ 是有限个开集的交而仍开。若 $z\in U\cap V$，选 $j$ 使 $z\in V_{y_j}$，又有 $z\in U_{y_j}$，与二者不交矛盾。若 $F=\varnothing$，直接取 $U=X,V=\varnothing$，亦与空交、空并约定相符。

不交是要求 $x\notin F$ 的直接原因。若去掉紧性，在 $\mathbb R$ 中取 $F=(0,1)$、$x=0$，任何 $0$ 的开邻域都与 $F$ 相交，结论失败。若去掉 Hausdorff，两点平凡拓扑空间中的单点 $F$ 已紧，却无法与另一点分离。
\end{remark}
% END SUPPLEMENT: id=v14-ch04-point-compact-details
% SOURCE p=143; printed=128; object=ch04-compact-subset-closed; number=4.24; region=bottom
\begin{proposition}\label{proposition:ch04-compact-subset-closed}
Hausdorff 空间中的每个紧子集都闭。
\end{proposition}
% SOURCE p=143; printed=128; object=ch04-proof-compact-subset-closed; region=bottom
\begin{proof}
由\cref{proposition:ch04-point-compact-separation}，若 $F$ 紧，则 $F^c$ 是其每个点的邻域，故 $F^c$ 开。
\end{proof}
% SUPPLEMENT: id=v14-ch04-compact-closed-details; parent=ch04-compact-subset-closed; kind=proof-detail+motivation+hypotheses; origin=independent-derivation
\begin{remark}
\textbf{动机与逐条件说明（重构新增）。}对每个 $x\in F^c$，前一命题给开集 $U_x\ni x$，与包含 $F$ 的 $V_x$ 不交，故 $U_x\subseteq F^c$。于是 $F^c=\bigcup_{x\in F^c}U_x$ 开，明确补全了“是每个点的邻域”的一步。紧性控制点不能在集合外形成未被捕获的极限，Hausdorff 性则提供排除这些外点的开邻域。两条件的作用和删去后的反例均见前一说明；例如两点平凡空间的真单点集紧但不闭。
\end{remark}
% END SUPPLEMENT: id=v14-ch04-compact-closed-details
% SOURCE p=144; printed=129; object=ch04-compact-intersections-unions; region=top
\begin{remark}
在非 Hausdorff 空间中，紧集未必闭（例如平凡拓扑空间的每个子集都紧），而紧集的交也未必紧，见\cref{exercise:ch04-doubled-origin}。当然，在 Hausdorff 空间中，由\cref{proposition:ch04-closed-subset-compact,proposition:ch04-compact-subset-closed}，任意紧集族的交都紧。此外，在任意拓扑空间中，有限个紧集的并总是紧的。事实上，若 $K_1,\ldots,K_n$ 紧，且 $\{U_\alpha\}$ 是 $\bigcup_{j=1}^nK_j$ 的开覆盖，只需对每个 $K_j$ 选有限子覆盖，再将它们合并。
\end{remark}
% SUPPLEMENT: id=v14-ch04-compact-set-operations; parent=ch04-compact-intersections-unions; kind=proof+correction+motivation+hypotheses; origin=independent-derivation
\begin{remark}
\textbf{空族的澄清（重构新增）。}原文“任意紧集族的交”在一般 Hausdorff 环境中须理解为\emph{非空族}。空族在 $X$ 中的交是 $X$，而 $X$ 未必紧，例如 $X=\mathbb R$。若环境本身也紧，则允许空族。以下按这一精确表述证明。
\end{remark}
\begin{proof}
\textbf{集合运算的完整补写（重构新增）。}平凡拓扑的任一子空间也只有空集与全空间两个开集，所以其任意开覆盖只要覆盖非空空间，就必含全空间，因而有单元素子覆盖；空子空间也紧。

若 $\{K_i\}_{i\in I}$ 是 Hausdorff 空间中的非空紧集族，每个 $K_i$ 闭，故 $K=\bigcap_iK_i$ 闭。选 $i_0\in I$，则 $K$ 是紧空间 $K_{i_0}$ 的相对闭子集，因而紧。总交本身可以为空，不影响论证。

有限并的证明中，若覆盖第 $j$ 个紧集所需的有限指标集为 $A_j$，则 $A_1\cup\cdots\cup A_n$ 仍有限且覆盖整个并。$n=0$ 时并为空，仍紧。无限并一般不保紧，例如每个 $[-n,n]$ 紧而其并 $\mathbb R$ 不紧。交的证明需要 Hausdorff 来获得闭性，还需要非空族来选一个紧的容纳空间；有限并的证明不需要任何分离公理。
\end{proof}
% END SUPPLEMENT: id=v14-ch04-compact-set-operations
% SOURCE p=144; printed=129; object=ch04-compact-hausdorff-normal; number=4.25; region=upper
\begin{proposition}\label{proposition:ch04-compact-hausdorff-normal}
每个紧 Hausdorff 空间都是正规的。
\end{proposition}
% SOURCE p=144; printed=129; object=ch04-proof-compact-hausdorff-normal; region=upper
\begin{proof}
设 $X$ 紧且 Hausdorff，$E,F$ 是其不交闭子集。由\cref{proposition:ch04-point-compact-separation}，对每个 $x\in E$，存在不交开集 $U_x,V_x$，使 $x\in U_x$、$F\subseteq V_x$。由\cref{proposition:ch04-closed-subset-compact}，$E$ 紧，而 $\{U_x\}_{x\in E}$ 覆盖 $E$，故可选有限子覆盖 $\{U_{x_j}\}_{j=1}^n$。令 $U=\bigcup_{j=1}^nU_{x_j}$、$V=\bigcap_{j=1}^nV_{x_j}$，则 $U,V$ 是不交开集，且 $E\subseteq U,F\subseteq V$。
\end{proof}
% SUPPLEMENT: id=v14-ch04-compact-normal-details; parent=ch04-compact-hausdorff-normal; kind=proof-detail+motivation+hypotheses; origin=independent-derivation
\begin{remark}
\textbf{动机与逐条件说明（重构新增）。}先用 $F$ 的闭性与 $X$ 的紧性得 $F$ 紧，才可在原证明第一步使用点与紧集分离；再用 $E$ 的闭性与紧性有限化其覆盖。若 $z\in U\cap V$，则某个 $j$ 满足 $z\in U_{x_j}$，而 $z\in V\subseteq V_{x_j}$，矛盾。空闭集情形可直接取 $\varnothing$ 和 $X$ 为分离开集。本书的“正规”还包括 $T_1$；这由 Hausdorff 已经保证，不能只证明闭集分离而漏掉该约定。此结果把两次有限化连接起来：第一次从点对分离到点与紧集分离，第二次到两个闭集分离。紧性和 Hausdorff 是这里的充分条件，并不声称每个正规空间都紧，例如通常实直线正规但不紧。
\end{remark}
% END SUPPLEMENT: id=v14-ch04-compact-normal-details
% SOURCE p=144; printed=129; object=ch04-continuous-image-compact; number=4.26; region=middle
\begin{proposition}\label{proposition:ch04-continuous-image-compact}
若 $X$ 紧且 $f:X\to Y$ 连续，则 $f(X)$ 紧。
\end{proposition}
% SOURCE p=144; printed=129; object=ch04-proof-continuous-image-compact; region=middle
\begin{proof}
设 $\{V_\alpha\}$ 是 $Y$ 中覆盖 $f(X)$ 的开集族。则 $\{f^{-1}(V_\alpha)\}$ 是 $X$ 的开覆盖，故有有限子覆盖 $\{f^{-1}(V_{\alpha_j})\}_{j=1}^n$。于是 $\{V_{\alpha_j}\}_{j=1}^n$ 就是 $f(X)$ 的有限子覆盖。
\end{proof}
% SUPPLEMENT: id=v14-ch04-compact-image-remark; parent=ch04-continuous-image-compact; kind=proof-detail+motivation+hypotheses; origin=independent-derivation
\begin{remark}
\textbf{动机与逐条件说明（重构新增）。}连续性使开覆盖的逆像仍开，紧性使逆像覆盖可有限化，而 $y\in f(X)$ 可写成 $f(x)$，从而该有限族也覆盖每个像点。只讨论 $f(X)$，所以不需要 $f$ 满射到 $Y$；也不需要单射或任何分离公理。连续性不能普遍删去：令 $f(0)=0$、$f(t)=1/t$（$0<t\le1$），紧区间 $[0,1]$ 的像 $\{0\}\cup[1,\infty)$ 不紧。若没有定义域紧性，实直线的恒等映射也不给紧像。
\end{remark}
% END SUPPLEMENT: id=v14-ch04-compact-image-remark
% SOURCE p=144; printed=129; object=ch04-compact-continuous-bounded; number=4.27; region=middle
\begin{corollary}\label{corollary:ch04-compact-continuous-bounded}
若 $X$ 紧，则 $C(X)=BC(X)$。
\end{corollary}
% SUPPLEMENT: id=v14-ch04-compact-bounded-proof; parent=ch04-compact-continuous-bounded; kind=proof+motivation+hypotheses; origin=independent-derivation
\begin{proof}
\textbf{原书省略证明的补写（重构新增）。}总有 $BC(X)\subseteq C(X)$。对 $f\in C(X)$，由上一命题，$f(X)$ 是 $\mathbb C$ 的紧子集。开圆盘 $B(0,n)$（$n\in\mathbb N$）覆盖 $\mathbb C$，故有限多个便覆盖 $f(X)$；取这些整数的最大值 $N$，得到 $|f(x)|<N$ 对所有 $x\in X$ 成立。所以 $f\in BC(X)$。若 $X$ 为空，则唯一空函数自动有界，等式仍成立。
\end{proof}
\begin{remark}
\textbf{动机与条件说明（重构新增）。}每点附近有界的连续函数，在紧空间上可把局部界压缩成一个全局界。这里 $C(X)$ 的值域约定为 $\mathbb C$，使用的是其通常度量及由整数半径圆盘组成的递增覆盖；实值情形同理。不要求 $X$ Hausdorff。紧性不能从这个等式反推，一般反向需要进一步条件，并且即使加正规也只得到后面习题中的可数紧性。
\end{remark}
% END SUPPLEMENT: id=v14-ch04-compact-bounded-proof
% SOURCE p=144; printed=129; object=ch04-compact-bijection-homeomorphism; number=4.28; region=middle-lower
\begin{proposition}\label{proposition:ch04-compact-bijection-homeomorphism}
若 $X$ 紧且 $Y$ Hausdorff，则任意连续双射 $f:X\to Y$ 都是同胚。
\end{proposition}
% SOURCE p=144; printed=129; object=ch04-proof-compact-bijection-homeomorphism; region=middle-lower
\begin{proof}
若 $E\subseteq X$ 闭，则依次由\cref{proposition:ch04-closed-subset-compact,proposition:ch04-continuous-image-compact,proposition:ch04-compact-subset-closed}，$E$ 紧，$f(E)$ 紧，且 $f(E)$ 闭。这表示 $f^{-1}$ 连续，故 $f$ 为同胚。
\end{proof}
% SUPPLEMENT: id=v14-ch04-compact-bijection-details; parent=ch04-compact-bijection-homeomorphism; kind=proof-detail+motivation+hypotheses; origin=independent-derivation
\begin{remark}
\textbf{动机与逐条件说明（重构新增）。}双射使反函数 $g=f^{-1}$ 存在；对 $X$ 的闭集 $E$，$g^{-1}(E)=f(E)$，所以证明 $f$ 是闭映射恰好证明 $g$ 连续。定义域紧性使闭集紧，$f$ 连续使像紧，值域 Hausdorff 使紧像闭，每一条件各负责一环。不另要求 $X$ Hausdorff，因为结论成立后它会从同胚自动获得该性质。

若删去紧性，从离散拓扑的 $\mathbb R$ 到通常拓扑的 $\mathbb R$ 的恒等双射连续，但反函数不连续；若删去值域 Hausdorff，从两点离散空间到两点平凡拓扑空间的恒等双射同样是连续双射而非同胚。双射不是装饰条件：没有它根本不存在两侧互逆映射。
\end{remark}
% END SUPPLEMENT: id=v14-ch04-compact-bijection-details
% SOURCE p=144; printed=129; object=ch04-compact-nets-intro; region=lower
现在证明紧拓扑空间具有一种 Bolzano--Weierstrass 性质。正如可以猜到的，只需把序列替换为网。
% SOURCE p=144; printed=129; object=ch04-compact-nets; number=4.29; region=lower
\begin{theorem}\label{theorem:ch04-compact-nets}
对拓扑空间 $X$，下列条件等价：
\begin{enumerate}
\item\label{item:ch04-compact-nets-cover} $X$ 紧；
\item\label{item:ch04-compact-nets-cluster} $X$ 中的每个网都有簇点；
\item\label{item:ch04-compact-nets-subnet} $X$ 中的每个网都有收敛子网。
\end{enumerate}
\end{theorem}
% SOURCE p=144,145; printed=129,130; object=ch04-proof-compact-nets; region=bottom-top
\begin{proof}
后两项的等价性由\cref{proposition:ch04-cluster-subnet}得到。若 $X$ 紧且 $\langle x_\alpha\rangle_{\alpha\in A}$ 是 $X$ 中的网，令 $E_\alpha=\{x_\beta:\beta\gtrsim\alpha\}$。因为对任意 $\alpha,\beta\in A$ 都存在 $\gamma\in A$，使 $\gamma\gtrsim\alpha,\beta$，集合族 $\{E_\alpha\}_{\alpha\in A}$ 有有限交性质。因此由\cref{proposition:ch04-compact-fip}，$\bigcap_{\alpha\in A}\overline{E_\alpha}\ne\varnothing$。若 $x$ 属于此交且 $U$ 为 $x$ 的邻域，则 $U$ 遇到每个 $E_\alpha$，也就是说此网频繁进入 $U$，所以 $x$ 为簇点。

另一方面，若 $X$ 不紧，取一个没有有限子覆盖的开覆盖 $\{U_\beta\}_{\beta\in B}$。令 $\mathcal A$ 为 $B$ 的有限子集族，按包含关系定向；对每个 $A\in\mathcal A$，选 $x_A\in(\bigcup_{\beta\in A}U_\beta)^c$。则 $\langle x_A\rangle_{A\in\mathcal A}$ 是一个没有簇点的网。事实上，给定 $x\in X$，选 $\beta\in B$ 使 $x\in U_\beta$。若 $A\supseteq\{\beta\}$，则 $x_A\notin U_\beta$，所以 $x$ 不是簇点。
\end{proof}
% SUPPLEMENT: id=v14-ch04-compact-nets-details; parent=ch04-compact-nets; kind=proof-detail+motivation+hypotheses; origin=independent-derivation
\begin{proof}
\textbf{有限交及反证构造的补写（重构新增）。}两个指标有共同上界，归纳后任意有限多个指标 $\alpha_1,\ldots,\alpha_n$ 也有共同上界 $\gamma$。由自反性及传递性，$x_\gamma\in\bigcap_{j=1}^nE_{\alpha_j}$，故该交非空。空子族的交为 $X$；存在一个以非空有向集为指标的 $X$ 值网已保证 $X\ne\varnothing$。于是尾集族以及闭包族均有有限交性质，才能对后者应用闭集刻画。

逆向中，每个有限并都不等于 $X$，保证选点集合非空；逐有限集选点使用选择公理。$\mathcal A$ 含空集，两个成员的并是共同上界，所以确为非空有向集。对任一 $x$，原证明给出其一个开邻域 $U_\beta$，且整个阈值 $\{\beta\}$ 之后的网值都避开它，违反频繁进入的条件。因此所有点都被排除，所得网无簇点。若 $X=\varnothing$，紧性成立，而另两项因不存在以非空指标集取值于 $X$ 的网而自动成立，等价性仍然有效。
\end{proof}
\begin{remark}
\textbf{动机与逐条件说明（重构新增）。}尾部闭包记录“无论走多远仍可能靠近”的点，紧性用有限交性质保证这些候选点不会全部消失。逆向的指标遍历所有有限覆盖片段，网值始终逃离已经检查过的片段，正把“不可能有限覆盖”转成“没有簇点”。不要求 Hausdorff，因而也不要求子网极限唯一；必须让网遍历任意非空有向指标集，不能在一般拓扑空间中只测试序列。
\end{remark}
% END SUPPLEMENT: id=v14-ch04-compact-nets-details
% SOURCE p=145; printed=130; object=ch04-countable-compact-intro; region=upper-middle
最后介绍另外两个与紧性有关的有用概念。
% SOURCE p=145; printed=130; object=ch04-countable-sequential-compactness; region=upper-middle
\begin{definition}[可数紧与序列紧]\label{definition:ch04-countable-sequential-compactness}
若拓扑空间 $X$ 的每个可数开覆盖都有有限子覆盖，则称 $X$ \textbf{可数紧}。若 $X$ 中的每个序列都有收敛子序列，则称 $X$ \textbf{序列紧}。
\end{definition}
% SOURCE p=145; printed=130; object=ch04-compactness-comparison; region=middle
当然，每个紧空间都可数紧；在度量空间中，紧性与序列紧性等价。然而在一般情形下，紧性与序列紧性互不蕴含。更多结果与例子见\cref{exercise:ch04-sequential-countably-compact,exercise:ch04-countably-compact-cluster,exercise:ch04-countably-compact-accumulation,exercise:ch04-ordinal-compactness,exercise:ch04-binary-digit-sequence}。
% SUPPLEMENT: id=v14-ch04-countable-compact-remark; parent=ch04-countable-sequential-compactness; kind=proof-detail+motivation+hypotheses; origin=independent-derivation
\begin{remark}
\textbf{动机与逐条件说明（重构新增）。}可数紧性限制覆盖族的大小，序列紧性限制探测收敛的指标集为 $\mathbb N$；两者是不同方向的可数化。这里“可数”包含有限，有限开覆盖本身已是有限子覆盖；“收敛子序列”要求极限在 $X$ 内，不能只在更大的环境中收敛。空空间两者都满足。紧性蕴含可数紧性，是因为可数开覆盖只是所有开覆盖中的一部分；度量空间的等价性已由\cref{theorem:ch00-compactness-characterization}证明。下述序数例子和二进制位例子将分别给出序列紧而非紧、紧而非序列紧的空间；第二例的紧性使用后续乘积定理，因此该部分依赖将在那里完成。
\end{remark}
% END SUPPLEMENT: id=v14-ch04-countable-compact-remark
\Needspace{7\baselineskip}
% SOURCE p=145; printed=130; object=ch04-compact-exercises-title; region=middle
\subsection*{习题}
% SOURCE p=145; printed=130; object=ch04-ex37-doubled-origin; number=37; region=middle
\begin{exercise}\label{exercise:ch04-doubled-origin}
设 $0'$ 是不属于 $(-1,1)$ 的一点，令 $X=(-1,1)\cup\{0'\}$。令 $\mathcal T$ 为下列集合生成的拓扑：
\begin{align*}
(-1,a),\quad(a,1),\quad
\bigl((-1,b)\setminus\{0\}\bigr)\cup\{0'\},\quad
\bigl((c,1)\setminus\{0\}\bigr)\cup\{0'\},
\end{align*}
其中 $-1<a<1$、$0<b<1$、$-1<c<0$。（可以把 $X$ 想成将 $(-1,1)$ 中的点 $0$ 分裂成两个点所得的空间。）
\begin{enumerate}
\item\label{item:ch04-doubled-origin-charts} 定义 $f,g:(-1,1)\to X$：对所有 $x$ 令 $f(x)=x$；对 $x\ne0$ 令 $g(x)=x$，并令 $g(0)=0'$。则 $f,g$ 都是到各自值域的同胚。
\item\label{item:ch04-doubled-origin-separation} $X$ 是 $T_1$ 空间但非 Hausdorff，尽管其中每点都有一个同胚于 $(-1,1)$ 的邻域（该邻域因而 Hausdorff）。
\item\label{item:ch04-doubled-origin-compacts} 集合 $[-\tfrac12,\tfrac12]$ 与 $([-\tfrac12,\tfrac12]\setminus\{0\})\cup\{0'\}$ 在 $X$ 中紧但不闭，而它们的交不紧。
\end{enumerate}
\end{exercise}
% SUPPLEMENT: id=v14-ch04-ex37-solution; parent=ch04-ex37-doubled-origin; kind=solution+motivation+hypotheses; origin=independent-derivation
\begin{solution}
\textbf{补写解答（重构新增）。}记 $I=(-1,1)$、$Y=X\setminus\{0'\}$、$Z=X\setminus\{0\}$。
\begin{enumerate}
\item $f:I\to Y$ 和 $g:I\to Z$ 都为双射。对题给子基，$f$ 的逆像是通常开射线或去掉 $0$ 的开射线，因而均开。$g$ 对普通射线的逆像是该射线去掉 $0$，对两种带 $0'$ 的射线的逆像则分别为 $(-1,b)$、$(c,1)$，也均开。因此两映射都连续。

通常开射线的 $f$-像是 $Y$ 中相对开集。通常下射线 $(-1,a)$ 的 $g$-像，在 $a>0$ 时为题给带 $0'$ 的下射线，在 $a\le0$ 时仍为普通下射线；通常上射线的情形对称。故这些像在 $Z$ 中相对开。射线的有限交与任意并组成 $I$ 中所有开集，双射保持这些运算，所以 $f,g$ 都是到各自值域的开映射。其逆映射因此连续，得到两同胚。
\item $Y$ 可由两个普通射线的并给出，$Z$ 可由两个带 $0'$ 的射线的并给出，所以二者在 $X$ 中都开。于是 $\{0'\}^c=Y$、$\{0\}^c=Z$ 开。若 $t\in I\setminus\{0\}$，则
\begin{align*}
X\setminus\{t\}=f(I\setminus\{t\})\cup g(I\setminus\{t\})
\end{align*}
也是开集：两像分别在开子空间 $Y,Z$ 中开。因此所有单点闭，$X$ 为 $T_1$。

任何含 $0$ 的开集包含一个含 $0$ 的有限子基交；该交只能使用普通射线，并含某个 $(-\varepsilon,\varepsilon)$。任何含 $0'$ 的有限子基交只能使用带 $0'$ 的射线，并含 $(( -\delta,\delta)\setminus\{0\})\cup\{0'\}$。这些邻域必在非零实数点相交，故 $0,0'$ 无法以不交开邻域分离，$X$ 非 Hausdorff。另一方面，$Y,Z$ 覆盖 $X$，本身都是同胚于 $I$ 的开邻域，证明局部断言。
\item 记题中两集合为 $K=f([-\tfrac12,\tfrac12])$、$L=g([-\tfrac12,\tfrac12])$。实闭区间紧，连续像也紧，所以 $K,L$ 紧。每个 $0'$ 的邻域都含非零且绝对值小于 $1/2$ 的点，故 $0'\in\overline K\setminus K$；同理 $0\in\overline L\setminus L$，所以二者都不闭。

两集合的交为 $H=[-\tfrac12,\tfrac12]\setminus\{0\}$，包含于 $Y$，其相对拓扑由第一项就是通常相对拓扑。集合
\begin{align*}
W_n=H\cap\bigl((-1,-1/n)\cup(1/n,1)\bigr),\qquad n\ge3,
\end{align*}
相对开且递增，并覆盖 $H$；任意有限多个只覆盖其中最大的 $W_N$，却漏掉 $0<|x|\le1/N$ 的点。因此 $H$ 不紧。
\end{enumerate}
\end{solution}
\begin{remark}
\textbf{反例的动机（重构新增）。}两个原点分别有通常区间的局部拓扑，却共享所有足够近的非零点，所以局部 Hausdorff 性无法保证全局 Hausdorff 性。这个同一构造同时说明：仅用 $T_1$ 不能替代“紧子集闭”及“紧集的非空交紧”中的 Hausdorff 条件。
\end{remark}
% END SUPPLEMENT: id=v14-ch04-ex37-solution
% SOURCE p=145; printed=130; object=ch04-ex38-compact-hausdorff-topologies; number=38; region=middle-lower
\begin{exercise}\label{exercise:ch04-compact-hausdorff-topologies}
设 $(X,\mathcal T)$ 是紧 Hausdorff 空间，$\mathcal T'$ 是 $X$ 上另一拓扑。若 $\mathcal T'$ 严格强于 $\mathcal T$，则 $(X,\mathcal T')$ 是 Hausdorff 但不紧；若 $\mathcal T'$ 严格弱于 $\mathcal T$，则 $(X,\mathcal T')$ 紧但非 Hausdorff。
\end{exercise}
% SUPPLEMENT: id=v14-ch04-ex38-solution; parent=ch04-ex38-compact-hausdorff-topologies; kind=solution+hypotheses; origin=independent-derivation
\begin{solution}
\textbf{补写解答（重构新增）。}若 $\mathcal T\subsetneq\mathcal T'$，原来分离不同点的不交开集仍在 $\mathcal T'$ 中，故较强拓扑仍 Hausdorff。恒等映射 $\operatorname{id}:(X,\mathcal T')\to(X,\mathcal T)$ 连续且双射。若较强拓扑也紧，由\cref{proposition:ch04-compact-bijection-homeomorphism}，此映射为同胚，迫使两拓扑相等，矛盾。因此较强拓扑不紧。

若 $\mathcal T'\subsetneq\mathcal T$，恒等映射 $\operatorname{id}:(X,\mathcal T)\to(X,\mathcal T')$ 连续，故其像即较弱拓扑空间紧。若该空间还 Hausdorff，同一命题再次给出同胚，迫使两拓扑相等，矛盾。因此它不 Hausdorff。严格包含用于排除相等；拓扑较强有利于分离却增加需检验的开覆盖，较弱有利于紧性却可能失去分离。
\end{solution}
% END SUPPLEMENT: id=v14-ch04-ex38-solution
% SOURCE p=145; printed=130; object=ch04-ex39-sequential-countably-compact; number=39; region=lower
\begin{exercise}\label{exercise:ch04-sequential-countably-compact}
每个序列紧空间都可数紧。
\end{exercise}
% SUPPLEMENT: id=v14-ch04-ex39-solution; parent=ch04-ex39-sequential-countably-compact; kind=solution; origin=independent-derivation
\begin{solution}
\textbf{补写解答（重构新增）。}反设序列紧空间 $X$ 有可数开覆盖 $\{U_j\}_{j\ge1}$ 而没有有限子覆盖。对每个 $n$，选 $x_n\notin\bigcup_{j=1}^nU_j$。序列紧性给出严格递增指标 $n_k$ 和 $x\in X$，使 $x_{n_k}\to x$。取 $j$ 使 $x\in U_j$。收敛使充分大的 $k$ 满足 $x_{n_k}\in U_j$，而 $n_k\ge k$，所以充分大时又有 $n_k\ge j$，按选点条件 $x_{n_k}\notin U_j$，矛盾。有限覆盖本就有有限子覆盖，故所有可数开覆盖的情况都已处理。
\end{solution}
% END SUPPLEMENT: id=v14-ch04-ex39-solution
% SOURCE p=145; printed=130; object=ch04-ex40-countably-compact-cluster; number=40; region=lower
\begin{exercise}\label{exercise:ch04-countably-compact-cluster}
若 $X$ 可数紧，则 $X$ 中每个序列都有簇点。若 $X$ 还第一可数，则 $X$ 序列紧。
\end{exercise}
% SUPPLEMENT: id=v14-ch04-ex40-solution; parent=ch04-ex40-countably-compact-cluster; kind=solution+hypotheses; origin=independent-derivation
\begin{solution}
\textbf{补写解答（重构新增）。}设 $\{x_n\}$ 是 $X$ 中的序列，令 $F_n=\overline{\{x_k:k\ge n\}}$。各 $F_n$ 非空、闭且递减。如果它们的交为空，则 $\{X\setminus F_n\}_{n\ge1}$ 是可数开覆盖。可数紧性给出有限子覆盖；因补集递增，这等价于某个 $X\setminus F_N=X$，与 $F_N$ 非空矛盾。故有 $x\in\bigcap_nF_n$。任何 $x$ 的邻域都遇到每个尾部，故 $x$ 是序列簇点。

若 $X$ 第一可数，取 $x$ 处可数邻域基 $\{V_j\}_{j\ge1}$，令 $W_j=\bigcap_{i=1}^jV_i$。归纳选 $n_j>n_{j-1}$ 使 $x_{n_j}\in W_j$：这由簇点对每个邻域频繁进入保证。对任意邻域 $U\ni x$，有 $V_m\subseteq U$，所以 $j\ge m$ 时 $x_{n_j}\in W_j\subseteq V_m\subseteq U$。因此 $x_{n_j}\to x$，得到序列紧性。这也逐步重现了\cref{exercise:ch04-sequence-cluster-subsequence}的结论。
\end{solution}
\begin{remark}
\textbf{条件的作用（重构新增）。}第一部分只用可数个尾部闭包，故可数紧已经足够；第二部分的第一可数性将所有邻域压缩成可递次满足的可数清单。这里不要求 $T_1$ 或 Hausdorff，因为序列可以反复取同一点，簇点定义允许这种重复。
\end{remark}
% END SUPPLEMENT: id=v14-ch04-ex40-solution
% SOURCE p=145; printed=130; object=ch04-ex41-countably-compact-accumulation; number=41; region=lower
\begin{exercise}\label{exercise:ch04-countably-compact-accumulation}
$T_1$ 空间 $X$ 可数紧，当且仅当 $X$ 的每个无限子集都有聚点。
\end{exercise}
% SUPPLEMENT: id=v14-ch04-ex41-solution; parent=ch04-ex41-countably-compact-accumulation; kind=solution+hypotheses; origin=independent-derivation
\begin{solution}
\textbf{补写解答（重构新增）。}先设 $X$ 可数紧，$E\subseteq X$ 无限。从中选互异点 $x_1,x_2,\ldots$。上一题给序列的簇点 $x$。对任意邻域 $U\ni x$，因序列各值互异，至多一项等于 $x$；从比该项更后的尾部选择落入 $U$ 的项，便得 $U\cap(E\setminus\{x\})\ne\varnothing$。因此 $x$ 是 $E$ 的聚点。这个方向甚至不需要 $T_1$。

反之，假设每个无限子集有聚点，却有无有限子覆盖的可数开覆盖 $\{U_j\}_{j\ge1}$。选 $x_n\notin\bigcup_{j=1}^nU_j$，令 $D=\{x_n:n\ge1\}$。$D$ 必无限：若有限，可为它的每个点选一个覆盖成员，再取这些成员指标的最大值 $N$，便得 $D\subseteq\bigcup_{j=1}^NU_j$，与 $x_N$ 的选择矛盾。

取 $D$ 的聚点 $x$，并取 $j$ 使 $x\in U_j$。由于 $n\ge j$ 时 $x_n\notin U_j$，有 $U_j\cap D\subseteq\{x_1,\ldots,x_{j-1}\}$。令 $F=(U_j\cap D)\setminus\{x\}$，这是有限集，在 $T_1$ 空间中闭。于是 $U_j\setminus F$ 是 $x$ 的开邻域，却不含 $D\setminus\{x\}$ 的点，与聚点定义矛盾。故 $X$ 可数紧。
\end{solution}
\begin{remark}
\textbf{$T_1$ 条件的作用与反例（重构新增）。}逆向需要从邻域中删去有限个其他点后仍是邻域。若删去 $T_1$，结论会失效：令 $X=\mathbb N\times\{0,1\}$，规定开集恰为各二点块 $B_n=\{n\}\times\{0,1\}$ 的任意并。覆盖 $\{B_n\}$ 无有限子覆盖，所以 $X$ 不可数紧。但任何非空子集 $E$，选其中一点 $(n,i)$，另一点 $(n,1-i)$ 的每个邻域都含 $(n,i)$，因此另一点是 $E$ 的聚点。特别是每个无限子集都有聚点。这说明缺少分离性时，“聚点”可能只来自无法分开的两点，而非无限多个不同点的趋近。
\end{remark}
% END SUPPLEMENT: id=v14-ch04-ex41-solution
% SOURCE p=145; printed=130; object=ch04-ex42-ordinal-compactness; number=42; region=lower
\begin{exercise}\label{exercise:ch04-ordinal-compactness}
可数序数集（\cref{definition:ch00-countable-ordinals-definition}）赋予序拓扑（\cref{exercise:ch04-order-topology}）后，序列紧且第一可数，但不紧。（证明序列紧性时使用\cref{proposition:ch00-countable-ordinals-bounded}。）
\end{exercise}
% SUPPLEMENT: id=v14-ch04-ex42-solution; parent=ch04-ex42-ordinal-compactness; kind=solution+motivation+hypotheses; origin=independent-derivation
\begin{solution}
\textbf{补写解答（重构新增）。}记可数序数集为 $\Omega$，最小元为 $o$；对 $\alpha\in\Omega$ 记 $I_\alpha=\{\xi:\xi<\alpha\}$、$K_\alpha=\{\xi:\xi\le\alpha\}$。由定义，每个 $I_\alpha$ 至多可数，而 $\Omega$ 不可数。因此 $\Omega$ 无最大元，否则它等于某个可数 $I_\alpha$ 加一点。良序性使每个 $\alpha$ 有直接后继 $\alpha^+=\min\{\xi:\xi>\alpha\}$。

\emph{第一可数性。}最小元 $o$ 的单点集 $\{o\}=I_{o^+}$ 开。若 $\alpha=\beta^+$，则 $\{\alpha\}=(\beta,\alpha^+)$ 开。对不是最小元的其他点 $\alpha$，集合族
\begin{align*}
\{(\beta,\alpha^+):\beta<\alpha\}
\end{align*}
可数且为 $\alpha$ 的邻域基。确实，每个包含 $\alpha$ 的基本开集是有限多个上下射线的交；取所有下界的最大值，并将上界缩小到 $\alpha^+$，就得到其中的一个上述邻域。若没有下界，任取 $\beta<\alpha$ 即可。因此每点都有可数邻域基。

\emph{有界闭初段的紧性。}固定 $\alpha\in\Omega$，取 $K_\alpha$ 的任意相对开覆盖。若存在 $\xi\le\alpha$，使 $K_\xi$ 不能由此覆盖的有限子族覆盖，取最小的这种 $\xi$。$\xi\ne o$，因为单点可由一个成员覆盖。选覆盖成员 $V\ni\xi$；由序拓扑的基本邻域形式，存在 $\beta<\xi$，使 $(\beta,\xi]\subseteq V$。这一点也适用于后继点：可取其直接前驱作为 $\beta$；若基本邻域没有下界，任取 $\beta<\xi$。由 $\xi$ 的最小性，$K_\beta$ 有有限子覆盖，加上 $V$ 即覆盖 $K_\xi$，矛盾。所以 $K_\alpha$ 紧。

\emph{序列紧性。}给任意序列 $\{x_n\}\subseteq\Omega$，其取值集合可数，由\cref{proposition:ch00-countable-ordinals-bounded}存在 $\alpha$ 使全部取值落在 $K_\alpha$ 内。刚才证明的紧性使 $K_\alpha$ 可数紧，且第一可数性由 $\Omega$ 限制到此子空间。由\cref{exercise:ch04-countably-compact-cluster}，该序列有在 $K_\alpha$ 中收敛的子序列。子空间中的收敛也给出在 $\Omega$ 中的收敛，因此 $\Omega$ 序列紧。

\emph{非紧性。}每个 $I_\alpha$ 在序拓扑中开，因每点都有更大的点，$\{I_\alpha:\alpha\in\Omega\}$ 覆盖 $\Omega$。有限个初始段的并等于其中最大的一个 $I_\alpha$，而它漏掉 $\alpha$，故没有有限子覆盖。因此 $\Omega$ 不紧。
\end{solution}
\begin{remark}
\textbf{反例的动机（重构新增）。}任意序列只涉及可数多个点，都会被圈在某个紧初段中，所以序列看不到空间向不可数方向不断延伸的整体。不可数开覆盖却能检出这种延伸。该例同时是第一可数、序列紧、可数紧而不紧的空间；可数紧性来自\cref{exercise:ch04-sequential-countably-compact}，不需要另作假设。
\end{remark}
% END SUPPLEMENT: id=v14-ch04-ex42-solution
% SOURCE p=145,146; printed=130,131; object=ch04-ex43-binary-digit-sequence; number=43; region=bottom-top
\begin{exercise}\label{exercise:ch04-binary-digit-sequence}
对 $x\in[0,1)$，令 $\sum_{n=1}^{\infty}a_n(x)2^{-n}$（$a_n(x)=0$ 或 $1$）为 $x$ 的二进制展开；若 $x$ 为二进有理数，则选充分大的 $n$ 上 $a_n(x)=0$ 的展开。则 $\{0,1\}^{[0,1)}$ 中的序列 $\langle a_n\rangle$ 没有逐点收敛子序列。（因此，由 $\{0,1\}$ 上离散拓扑所诱导的乘积拓扑下，$\{0,1\}^{[0,1)}$ 不序列紧；但它是紧的，这将在\cref{section:ch04-tychonoff}中证明。）
\end{exercise}
% SUPPLEMENT: id=v14-ch04-ex43-solution; parent=ch04-ex43-binary-digit-sequence; kind=solution+hypotheses; origin=independent-derivation
\begin{solution}
\textbf{补写解答（重构新增）。}任取严格递增的指标 $n_1<n_2<\cdots$。定义二进制位 $b_n$：当 $n=n_k$ 且 $k$ 为奇数时令 $b_n=1$，其他情形令 $b_n=0$；再令
\begin{align*}
x=\sum_{n=1}^{\infty}b_n2^{-n}.
\end{align*}
它至少有一个 $1$，故 $x>0$；也至少有一个 $0$，故 $x<\sum_{n=1}^{\infty}2^{-n}=1$。这些位中有无穷多个 $0$ 和无穷多个 $1$，所以此展开既非最终为 $0$，也非最终为 $1$。

为排除展开歧义，设两个不同的二进制位列给出同一实数，并设第一个不同位为 $m$。该位贡献差的绝对值为 $2^{-m}$，而后续贡献绝对值之和至多 $\sum_{n>m}2^{-n}=2^{-m}$。要完全抵消，后面每位必须取到同一方向的最大差，也就是一个位列在 $m$ 后全为 $0$，另一个全为 $1$。我们构造的位列不属于这种情形，所以它就是 $x$ 的唯一二进制展开，特别与题目约定的 $a_n(x)$ 相同。

于是 $a_{n_k}(x)=1$（$k$ 奇）而 $a_{n_k}(x)=0$（$k$ 偶），故此子序列在坐标 $x$ 上不收敛。每个候选子序列都能如此选出一个使它振荡的坐标，因此没有逐点收敛子序列。由\cref{exercise:ch04-weak-net-convergence}，乘积拓扑中的序列收敛恰为各坐标收敛，故该乘积空间不序列紧。

每个因子 $\{0,1\}$ 为有限空间，因而紧。全乘积的紧性将在\cref{section:ch04-tychonoff}由任意紧空间乘积定理得到；这里先完成无收敛子序列这一断言，后续定理不依赖此习题。
\end{solution}
\begin{remark}
\textbf{构造动机与条件说明（重构新增）。}给定任何候选指标子序列，都有一个坐标专门在这些指标上交替取值。二进有理数的展开约定使每个 $a_n$ 真正成为单值函数；构造无限多个 $0$ 和 $1$ 则主动避开歧义。乘积拓扑只要求逐坐标收敛，这里只需破坏一个坐标便足够，无须证明统一振荡。
\end{remark}
% END SUPPLEMENT: id=v14-ch04-ex43-solution
% SOURCE p=146; printed=131; object=ch04-ex44-countably-compact-image; number=44; region=upper
\begin{exercise}\label{exercise:ch04-countably-compact-image}
若 $X$ 可数紧且 $f:X\to Y$ 连续，则 $f(X)$ 可数紧。
\end{exercise}
% SUPPLEMENT: id=v14-ch04-ex44-solution; parent=ch04-ex44-countably-compact-image; kind=solution; origin=independent-derivation
\begin{solution}
\textbf{补写解答（重构新增）。}设 $\{V_n\}_{n\ge1}$ 是 $f(X)$ 的可数相对开覆盖，写 $V_n=f(X)\cap U_n$，其中 $U_n$ 在 $Y$ 中开。则 $\{f^{-1}(U_n)\}$ 是 $X$ 的可数开覆盖。可数紧性给出有限指标集 $F$，使 $X=\bigcup_{n\in F}f^{-1}(U_n)$。对任意 $y=f(x)\in f(X)$，某个 $n\in F$ 满足 $x\in f^{-1}(U_n)$，故 $y\in V_n$，所以这些 $V_n$ 覆盖 $f(X)$。有限或空覆盖情形直接适用同一有限化结论。连续性用于保持逆像开性，指标族仍可数，且不要求 $Y$ 有分离性质。
\end{solution}
% END SUPPLEMENT: id=v14-ch04-ex44-solution
% SOURCE p=146; printed=131; object=ch04-ex45-normal-pseudocompact; number=45; region=upper
\begin{exercise}\label{exercise:ch04-normal-countably-compact-bounded}
若 $X$ 正规，则 $X$ 可数紧，当且仅当 $C(X)=BC(X)$。（使用\cref{exercise:ch04-countably-compact-cluster,exercise:ch04-countably-compact-image}。若 $\{x_n\}$ 是 $X$ 中没有簇点的序列，则 $\{x_n:n\in\mathbb N\}$ 闭，可以应用\cref{corollary:ch04-tietze-unbounded}。）
\end{exercise}
% SUPPLEMENT: id=v14-ch04-ex45-solution; parent=ch04-ex45-normal-pseudocompact; kind=solution+hypotheses; origin=independent-derivation
\begin{solution}
\textbf{补写解答（重构新增）。}先设 $X$ 可数紧。对任意 $f\in C(X)$，上一题给出 $f(X)\subseteq\mathbb C$ 可数紧。圆盘族 $\{B(0,n):n\ge1\}$ 限制到 $f(X)$ 为可数开覆盖，故有限多个便覆盖它。取这些半径的最大值，得到 $f$ 有界。因此 $C(X)=BC(X)$；这个方向不需要正规性。

反之，设 $X$ 正规但不可数紧。取无有限子覆盖的可数开覆盖 $\{U_j\}_{j\ge1}$，并选 $x_n\notin\bigcup_{j=1}^nU_j$。此序列没有簇点：对任意 $x\in X$，某个 $U_j$ 含 $x$，但从第 $j$ 项开始所有项都避开 $U_j$。其取值集合 $D$ 无限，证明同\cref{exercise:ch04-countably-compact-accumulation}中的选点论证。将 $D$ 无重复地枚举为 $\{d_k:k\ge1\}$。

现在证明 $D$ 闭且离散。对每个 $x\in X$，存在开邻域 $V_x$，只含该序列有限多个项，因此 $V_x\cap D$ 为有限集。若 $x\notin D$，由正规空间的 $T_1$ 性可从 $V_x$ 中删去这有限集，得到不遇 $D$ 的开邻域，故 $D$ 闭。若 $x=d_k\in D$，从 $V_x$ 中删去 $(V_x\cap D)\setminus\{d_k\}$，得到与 $D$ 恰交于 $d_k$ 的开邻域，所以 $D$ 的相对拓扑离散。

定义 $h:D\to\mathbb R$ 为 $h(d_k)=k$。因 $D$ 离散，$h$ 连续。由\cref{corollary:ch04-tietze-unbounded}，闭子集 $D$ 上的 $h$ 可延拓为连续 $H:X\to\mathbb R$。将它也视为复值连续函数，则 $H\in C(X)$，却因 $H(d_k)=k$ 而无界。所以 $C(X)\ne BC(X)$，证明逆向的逆否命题。
\end{solution}
\begin{remark}
\textbf{动机与逐条件说明（重构新增）。}可数紧性阻止连续函数沿可数点列逃向无穷。反过来，一个无簇点序列提供闭离散点集，可在上面任意指定不断增大的函数值，再用 Tietze 延拓把局部反例变成全空间反例。正规性在延拓时使用，而本书正规定义中的 $T_1$ 在删去有限点时使用。必须先把重复值去掉再规定 $h(d_k)=k$，否则同一个点可能被赋予多个数值；原提示关于闭性的断言也需上述有限点删除论证。
\end{remark}
% END SUPPLEMENT: id=v14-ch04-ex45-solution
\end{document}
